\section*{Capítulo \ref{ch:b2:geological-hazards} --- Ameaças e riscos geológicos}

\begin{solution}{exo:b2:geological-hazards:1}
(a) $1 - e^{-30/100} = 0.259$. (b) $e^{-100/100} = e^{-1} = 0.368$: mais de um século em cada três
passa sem nenhuma. (c) $T = -50/\ln 0.95 = 975$ anos.
\end{solution}

\begin{solution}{exo:b2:geological-hazards:2}
(a) Cidade A: $(1/200) \times 5000 \times 0.02 = 0.5$ edifício por ano; cidade B: $(1/50) \times 800
\times 0.10 = 1.6$. (b) A cidade B, com 6.25 vezes menos edifícios expostos. (c) Sua ameaça é 4 vezes
mais frequente e sua vulnerabilidade 5 vezes maior: juntas, um fator 20, que supera a
\omterm{def:b2:geological-hazards:risk}{exposição} menor.
\end{solution}

\begin{solution}{exo:b2:geological-hazards:3}
(a) $\sqrt{9.81 \times 4000} = \qty{198}{m.s^{-1}}$; $\sqrt{9.81 \times 200} = \qty{44.3}{m.s^{-1}}$;
$\sqrt{9.81 \times 10} = \qty{9.90}{m.s^{-1}}$. (b) $\num{1.5e6}/198.1 = \qty{7570}{s}$, \qty{2.10}{h}.
(c) $800/3.6 = \qty{222.2}{m.s^{-1}}$, e $h = c^2/g = 222.2^2/9.81 = \qty{5030}{m}$.
\end{solution}

\begin{solution}{exo:b2:geological-hazards:4}
(a) $\tan\ang{33}/\tan\ang{25} = 0.6494/0.4663 = 1.39$. (b) $F = 1$ quando $\beta = \varphi' =
\ang{33}$. (c) Com $c' = 0$, $F = (\cos^2\beta - r_u)\tan\varphi'/(\sin\beta\cos\beta)$, e $\cos^2\beta
- r_u = \cos^2\beta\,(1 - \gamma_w/\gamma)$, logo $F = (1 - \gamma_w/\gamma)\tan\varphi'/\tan\beta$:
a água reduz o atrito à metade. $(1 - 9.81/20) \times 1.393 = 0.510 \times 1.393 = 0.710$: a
encosta rompe.
\end{solution}

\begin{solution}{exo:b2:geological-hazards:5}
(a) $696/50 = 13.9$ por ano. (b) $e^{-30/26.2} = 0.318$, contra 0.304 observado: próximo. O
valor observado é um pouco menor porque as réplicas e os dupletos agrupados acrescentam intervalos
curtos. (c) $1 - e^{-7/26.2} = 0.234$.
\end{solution}

\begin{solution}{exo:b2:geological-hazards:6}
(a) O termo de magnitude é $-0.105 + 0.229 \times 0.5 = 0.0095$. Para $d = 0$, 10, 30 e
\qty{100}{km}, $r = 5.57$, 11.45, 30.51 e \qty{100.2}{km}, $\log_{10}\mathrm{PGA} = -0.571$,
$-0.814$, $-1.145$ e $-1.547$, e a PGA mediana é 0.269, 0.153, 0.0715 e \qty{0.0284}{g}.
(b) $0.153 \times 10^{0.251} = 0.153 \times 1.78 = \qty{0.272}{g}$. (c) Um desvio-padrão
de cada lado: de $0.153/1.70 = \qty{0.0900}{g}$ a $0.153 \times 1.70 = \qty{0.260}{g}$.
\end{solution}

\begin{solution}{exo:b2:geological-hazards:7}
(a) Em ordem: 420 ($T = 13/1 = 13$ anos), 305 (6.5), 268 (4.33), 251 (3.25). (b) $305 + (420 -
305)\ln(10/6.5)/\ln(13/6.5) = 305 + 115 \times 0.621 = \qty{376}{m^3.s^{-1}}$. (c) O maior pico
é plotado em 13 anos; qualquer valor em 100 anos extrapola muito além do registro, com uma lei de
extremos cuja escolha e cujos parâmetros os 12 pontos mal restringem.
\end{solution}

\begin{solution}{exo:b2:geological-hazards:8}
(a) $0.6 \times 0.4 \times 0.5 \times 0.3 = 0.036$. (b) $P(\text{erupção}) = 0.6 \times 0.4 = 0.24$;
sem fluxo no vale, $0.24 - 0.036 = 0.204$. (c) $0.6 \times 0.7 \times 0.5 \times 0.3 =
0.063$.
\end{solution}

\begin{solution}{exo:b2:geological-hazards:9}
Aqui $\sin\beta\cos\beta = 0.4698$, $\cos^2\beta = 0.6710$ e $\tan\ang{30}/\tan\ang{35} =
0.8245$. (a) $F = c'/(\gamma z\sin\beta\cos\beta) + 0.8245 = 1$ dá $z_c = 8/(19 \times 0.4698
\times 0.1755) = \qty{5.11}{m}$. (b) $8/(19 \times 2 \times 0.4698) + 0.8245 = 0.448 + 0.825 =
1.27$. (c) $r_u = 9.81 \times 0.6710/19 = 0.3465$; $F = [8 + 38 \times (0.6710 - 0.3465) \times
0.5774]/(38 \times 0.4698) = 15.12/17.85 = 0.847$: a camada saturada rompe.
\end{solution}

\begin{solution}{exo:b2:geological-hazards:10}
(a) $\Pr(\text{nenhum em } [0, s+t] \mid \text{nenhum em } [0, s]) = e^{-(s+t)/T}/e^{-s/T} = e^{-t/T}$.
(b) Os 400 anos tranquilos não mudam nada: $1 - e^{-50/475} = 0.0999$, \qty{10}{\percent}. (c) Ele
contradiz a hipótese de que a taxa não depende do tempo decorrido: depois do rebote a
tensão precisa acumular-se de novo, de modo que a taxa é baixa logo após um terremoto e cresce depois. Uma zona
sismogênica contém muitas falhas em fases diferentes de seus ciclos; a soma de muitas sequências independentes
de eventos comporta-se quase como um processo de Poisson, e a taxa é tudo o que se conhece de cada uma.
\end{solution}

\begin{solution}{exo:b2:geological-hazards:11}
(a) $A = 0.5 \times (4000/10)^{1/4} = 0.5 \times 4.47 = \qty{2.24}{m}$. (b) $\lambda = cT$:
$198.1 \times 1200 = \qty{238}{km}$ no oceano profundo, $9.90 \times 1200 = \qty{11.9}{km}$ em
\qty{10}{m} de água. (c) A amplitude já não é pequena diante da profundidade (\qty{2.2}{m} em
\qty{10}{m}), de modo que a onda se empina e pode quebrar; parte da energia é refletida ou concentrada pela
forma do fundo do mar e das baías; em terra a altura máxima de inundação depende da declividade e da rugosidade do
terreno.
\end{solution}

\begin{solution}{exo:b2:geological-hazards:12}
(a) $10/5.8 + 4 = 1.72 + 4 = \qty{5.72}{s}$. (b) A onda S atinge \qty{50}{km} depois de $50/3.2 =
\qty{15.6}{s}$: $15.6 - 5.72 = \qty{9.90}{s}$ de aviso; a \qty{100}{km}, $31.3 - 5.72 =
\qty{25.5}{s}$. (c) $3.2 \times 5.72 = \qty{18.3}{km}$.
\end{solution}

\begin{solution}{pb:b2:geological-hazards:1}
Os logaritmos são mantidos com quatro casas decimais nas etapas intermediárias.
\textbf{1.} $10^{4-5} = 0.1$, $10^{4-6} = 0.01$ e $10^{4-7} = 0.001$ por ano.
\textbf{2.} 10, 100 e 1000 anos.
\textbf{3.} $10^{4-7.5} = \num{3.16e-4}$ por ano; $0.001 - \num{3.16e-4} = \num{6.84e-4}$ por ano.
\textbf{4.} $0.1 - 0.01 = 0.09$ e $0.01 - 0.001 = 0.009$ por ano.
\textbf{5.} $1 - e^{-0.01 \times 50} = 0.393$.
\textbf{6.} $1 - e^{-0.001 \times 50} = 0.0488$.
\textbf{7.} $r = \sqrt{10^2 + 5.57^2} = \qty{11.45}{km}$; $\log_{10} r = 1.0588$ e $-0.778 \times
1.0588 = -0.8237$.
\textbf{8.} $\mu = -0.105 + 0.229(M - 6) - 0.8237$: $-1.0432$, $-0.8142$ e $-0.6424$; medianas
0.0905, 0.153 e \qty{0.228}{g}.
\textbf{9.} $10^{-0.8142 + 0.230} = \qty{0.260}{g}$.
\textbf{10.} $z = (\log_{10} 0.3 + 0.8142)/0.230 = (-0.5229 + 0.8142)/0.230 = 1.266$, e $1 -
\Phi(1.266) = 0.103$.
\textbf{11.} $-0.105 + 0.229 \times 1.5 - 0.8237 = -0.5852$: \qty{0.260}{g}.
\textbf{12.} As probabilidades de exceder \qty{0.2}{g} ($\log_{10} 0.2 = -0.6990$) são 0.0672
($z = 1.497$), 0.308 ($z = 0.501$) e 0.597 ($z = -0.246$): $\lambda = 0.09 \times 0.0672 + 0.009
\times 0.308 + \num{6.84e-4} \times 0.597 = \num{6.05e-3} + \num{2.77e-3} + \num{4.08e-4} =
\num{9.23e-3}$ por ano.
\textbf{13.} Para \qty{0.3}{g}: $z = 2.262$, 1.266 e 0.520, probabilidades 0.0118, 0.103 e
0.302; $\lambda = \num{1.07e-3} + \num{9.24e-4} + \num{2.06e-4} = \num{2.20e-3}$ por ano.
\textbf{14.} Para \qty{0.4}{g} ($\log_{10} 0.4 = -0.3979$): $z = 2.805$, 1.810 e 1.063,
probabilidades 0.00251, 0.0352 e 0.144; $\lambda = \num{2.26e-4} + \num{3.17e-4} + \num{9.84e-5} =
\num{6.41e-4}$ por ano.
\textbf{15.} $-\ln 0.9/50 = \num{2.107e-3}$ por ano, um \omterm{def:b2:geological-hazards:return}{período de retorno} de 475 anos.
\textbf{16.} Entre \qty{0.3}{g} (\num{2.20e-3}) e \qty{0.4}{g} (\num{6.41e-4}).
\textbf{17.} A inclinação é $\ln(6.41/22.0)/\ln(4/3) = -1.233/0.2877 = -4.29$, e $a = 0.3 \times
(2.107/2.20)^{-1/4.29} = 0.3 \times 1.0101 = \qty{0.303}{g}$.
\textbf{18.} A \qty{0.3}{g}, a classe 5--6 dá \qty{48.5}{\percent} da taxa, a 6--7
\qty{42.1}{\percent} e a 7--7.5 apenas \qty{9.4}{\percent}. Mesmo o maior terremoto tem uma mediana de
\qty{0.260}{g}, abaixo de \qty{0.3}{g}: o nível só é atingido na cauda superior da dispersão, e
são os terremotos moderados frequentes que fornecem a maior parte dele.
\textbf{19.} $4/(19 \times 3 \times 0.4330) + \tan\ang{34}/\tan\ang{30} = 0.162 + 1.168 = 1.33$.
\textbf{20.} $r_u = 9.81 \times 0.75/19 = 0.387$; $F = [4 + 57 \times (0.75 - 0.387) \times
0.6745]/(57 \times 0.4330) = 17.95/24.68 = 0.727$: a encosta saturada rompe.
\textbf{21.} Com o nível freático a uma fração $m$ da espessura, $r_u = 0.387\,m$. $F = 1$ quando
$0.75 - 0.387\,m = (24.68 - 4)/(57 \times 0.6745) = 0.538$: $m = 0.548$, um nível freático a
\qty{55}{\percent} da camada.
\textbf{22.} Dividindo $F = 1$ por $\gamma z$: $0.0702 + (0.75 - 0.4330\,k) \times 0.6745 = 0.4330 +
0.75\,k$, logo $k_c = (0.0702 + 0.5059 - 0.4330)/(0.2921 + 0.75) = 0.1431/1.042 = 0.137$.
\textbf{23.} $\log_{10} 0.137 = -0.8633 = -0.105 + 0.229(M - 6) - 0.8237$ dá $M = 6.29$. Esses
terremotos ocorrem a $10^{4 - 6.29} - 10^{-3.5} = \num{4.81e-3}$ por ano, cerca de uma vez em 208 anos,
e a PGA de 475 anos da questão~17 é mais que o dobro dos \qty{0.137}{g} críticos: segundo esta
análise, a encosta seca rompe com um abalo muito inferior ao movimento de 475 anos.
\textbf{24.} A \omterm{def:b2:geological-hazards:psha}{aceleração de pico do solo} com uma chance de \qty{10}{\percent} de ser excedida em 50
anos na cidade é de \textbf{\qty[text-series-to-math=true]{0.303}{g}}.
\end{solution}
